\documentclass{ximera}

\input{../preamble.tex}


\title[Problem 3]{Problem 3}

\begin{document}
\begin{abstract} \end{abstract}
\maketitle


% Extracted from indeterminateForms.tex, problem #3
\begin{problem}
Evaluate the following limits.
	\begin{enumerate}
		\item $ \displaystyle \lim{x \to \frac{\pi}{2}} \frac{\cos{(x)}}{x}  $
	\begin{explanation} Apply the Quotient Law and use continuity of the numerator and denominator.\\[1em]
	 $ \displaystyle \lim{x \to \frac{\pi}{2}} \frac{\cos{(x)}}{x}=\frac{0}{\frac{\pi}{2}}=0  $
	\end{explanation} 
		
	\item	$\displaystyle \lim{x \to 3} \frac{\sqrt{25-x^2 }-4}{x-4} $\\
	\begin{explanation}Apply the Quotient Law and use continuity.	\\[1em]
	 $ \displaystyle \lim{x \to 3} \frac{\sqrt{25-x^2 }-4}{x-4}=  \frac{\sqrt{25-(3)^2 }-4}{3-4}= \frac{4-4}{-1}=0$\\
	
	\end{explanation} 
	\item	$\displaystyle \lim{x \to 3} \frac{\sqrt{25-x^2 }-4}{x-3} $\\
	\begin{explanation}
	Since $\displaystyle \lim{x \to 3} (\sqrt{25-x^2 }-4) = 0$ and $\displaystyle \lim{x \to 3} {x-3} = 0$, this limit has form $\frac{0}{0}$.\\[1em]
	 $ \displaystyle \lim{x \to 3} \frac{\sqrt{25-x^2 }-4}{x-3}=  \displaystyle \lim{x \to 3} \frac{\sqrt{25-x^2 }-4}{x-3}\frac{\sqrt{25-x^2 }+4}{\sqrt{25-x^2 }+4}=$\\
	 $ \displaystyle \lim{x \to 3} \frac{25-x^2-16}{(x-3)(\sqrt{25-x^2 }+4)}=  \displaystyle \lim{x \to 3} \frac{9-x^2}{(x-3)(\sqrt{25-x^2 }+4)}= \displaystyle \lim{x \to 3} \frac{(3+x)(3-x)}{(x-3)(\sqrt{25-x^2 }+4)}=\\   $
	 $\displaystyle \lim{x \to 3} \frac{-(3+x)}{\sqrt{25-x^2 }+4}=\frac{-6}{\sqrt{25-9 }+4}=\frac{-6}{\sqrt{16 }+4}=\frac{-6}{8}=\frac{-3}{4}$
	\end{explanation}
\end{enumerate}
\end{problem}



\end{document}
